\scanpage{48}
\handwritten{3}) Relations entre les deux théories (pour schémas abéliens et pour
\deleted{\unsure{c)|2)} Relations} groupes de BT). Supposons que $p$ soit localeme\deleted{\unsure{j}}nt nilpotent sur $S$. Travaillons \deleted{On suppose} à nouveau avec le site cristallin de Berthel\deleted{l}ot (pas le
nilpotent). Le résul\deleted{\unsure{lt}}tat essentiel sera, pour un s\deleted{\unsure{v}}chéma abélien variable\deleted{\unsure{,}}
$A$ sur $S$, un isomorphisme fonctoriel compatible aux changements de base
\[
\underline D^{\handwritten{*}}(A(\infty))\handwritten{\simeq} R^{1}f_{\mathrm{cris}\,\handwritten{*}}(\underline O_{A_{\mathrm{cris}}})\quad ,
\]
cet isomorphisme induisant, pour les valeurs des deux membres sur $S$, un
isomorphisme compatible aux filtrations \qquad et \qquad (compte tenu des
isomorphismes $\underline t_{A(\infty)}\handwritten{\simeq}\underline t_A$ et idoine pour $A^{*}$).\inserted{\handwritten{-- C'est en fait en postulant
un tel isomorphisme qu'on arrive par voie heuristique à une définition de $\mathbb D^{*}(G)$ pour un groupe de BT $G$ quelconque}}
% slanted: the insertion continues in the left margin after « définition ».

Ceci impliquera en particulier le fait que, pour un \deleted{voisinage} épaississement à p. div. loc. nilp. \deleted{de} $S'$ de $S$, il y a identité de la théorie
des prolongements infinitésimaux de $A$, et de $A(\infty)$. De eeci, par dévissage
sur l'ordre de nilpot\deleted{\unsure{t}}ence, on peut déduire par exemple immédi\deleted{\unsure{q}}atement le
théorème de Serre-Tate déjà annoncé plus haut\handwritten{:}\inserted{\handwritten{ méthode suivie par
Messing ds sa thèse.}}

\handwritten{[}On peut aussi définir une extension vectorielle universe\deleted{\unsure{e}}lle \deleted{$\underline E$} $E(G)$
d'un groupe de BT, et on trouvera une isomorphisme canonique etc
\[
E(A(\infty))\handwritten{\simeq} E(A)(\infty)
\]
donnant naissance, par le procédé habituel, à un isomorphisme plus
général de cristaux en \deleted{G}groupes
\[
\underline E(A\handwritten{(}\infty\handwritten{))}\handwritten{\simeq}\underline E(A)(\infty)\quad\handwritten{....}\qquad\handwritten{]}
\]

\deleted{4) Passage à la limite}

\deleted{\handwritten{4}) Problème\unsure{.|:} trouver le foncteur des F-cristaux filtrés vers \deleted{\unsure{modules}}
faisceaux $p$-adiques.}
% archive sheet number: 47

